% La revue de littérature présente l'état de l'art dans le domaine
% de recherche. Elle situe le travail par rapport aux travaux existants
% et identifie les lacunes que la recherche vise à combler.
\Chapter{LITERATURE REVIEW}\label{sec:RevLitt}

%%
%% SECTION 1
%%
\section{Historical Background}

Solving an optimization problem on a computer requires translating the mathematical formulation
into a data format that solvers can process. For linear programs, this means constructing the
constraint matrix $A$, the right-hand side vector $b$, and the objective coefficients $c$ that
define the \emph{standard form},
\begin{equation}
\begin{array}{ll}
\mbox{minimize} & c^T x \\
\mbox{subject to} & Ax \leq b, \\
& x \geq 0.
\end{array}
\end{equation}

\paragraph{Manual encoding}
In the early days of linear programming, analysts manually computed and encoded each coefficient
onto punch cards. This process was tedious, error-prone and became impractical as linear
programming found more and more applications in large-scale industrial
problems~\cite{fourer1978modern}.

\paragraph{MPS format}
IBM's Mathematical Programming System (MPS) file format, developed in the 1960s, established a
standardized file format for representing linear programs~\cite{ibm1976mpsx}. The MPS format uses
fixed-width fields designed for 80-column punch cards, with sections for rows, columns, and
coefficients. 
William Orchard-Hays~\cite{orchard-hays1984history} provides a comprehensive history of these
early mathematical programming systems, tracing their development from 1952 at RAND Corporation
through their widespread adoption on mainframe computers in the late 1960s.

\paragraph{Matrix generators}
To automate the creation of MPS files, \emph{matrix generators} emerged in the late 1960s; they
are procedural programs that compute the problem coefficients through explicit loops and output
the results in MPS format~\cite{marceau1971matrix}. Commercial systems such as
\texttt{MaGen}~\cite{haverly1977magen}, \texttt{GAMMA}~\cite{sperry1977gamma}, and \texttt{DATAFORM}~\cite{ketron1975dataform}
were among the first matrix generators which were used in various domains.

\paragraph{A modern approach}
Modelers naturally think in terms of sets, indices, and algebraic expressions, not in terms of
individual matrix entries. Matrix generators simplified the encoding process but still required
modelers to translate their mathematical formulations into procedural code that explicitly
constructs the coefficient matrices. This suggests that the gap between the modeler's
representation of the problem and the solver's required format hadn't been resolved
yet~\cite{fourer1978modern}.

\paragraph{Algebraic modeling languages}
In the late 1970s, \emph{algebraic modeling languages} emerged as formal languages designed to
accept problem descriptions that closely mirror standard mathematical
notation~\cite{fourer1983modeling}. By working directly with the modeler's representation, these
languages simplify verification (the model can be inspected as written), modification (changing
the problem does not require rewriting procedural code), and communication (the formulation serves
as its own documentation). In the late 1980s and early 1990s, commercial algebraic modeling
systems such as \texttt{GAMS}~\cite{brooke1988gams}, \texttt{AIMMS}~\cite{bisschop2006aimms} and
\texttt{AMPL}~\cite{fourer1990ampl} became widely adopted in both academia and industry (and they are still
widely used today).
%%
%% SECTION 2
%%
\section{Related work}

In the 2000s, advances in convex optimization made it practical to solve a much broader class of
problems efficiently. Solvers such as \texttt{MOSEK} \cite{mosekSolver, aps2025mosek} adopted \emph{conic
optimization} as their interface---a structured subclass of convex optimization with the standard
form~\cite{boyd2004convex}:
\begin{equation}
      \begin{array}{ll}
        \text{minimize} & c^T x\\
        \text{subject to} &  b-Ax \in \mathcal K,\\
      \end{array}
      \label{eq:conic}
\end{equation}
where $\mathcal{K}$ is a closed convex cone. This formulation generalizes linear programming
(when $\mathcal{K}$ is the nonnegative orthant) to include second-order cone programs
(SOCPs), semidefinite programs (SDPs) and more.

\paragraph{Embedded modeling languages}
In the late 2000s and early 2010s, modeling languages emerged that are embedded in general-purpose
programming languages (in contrast with standalone systems like \texttt{AMPL} and \texttt{GAMS}). Examples include
\texttt{YALMIP}~\cite{lofberg2004yalmip} in MATLAB, \texttt{Pyomo}~\cite{hart2011pyomo} in Python,
\texttt{JuMP}~\cite{dunning2017jump} in Julia, and \texttt{CasADi}~\cite{andersson2019} in C++.

\paragraph{Disciplined convex programming}
In addition, there are also \emph{domain specific languages} that implement
\emph{disciplined convex programming} (DCP), a set of composition rules that can automatically
verify convexity~\cite{grant2004disciplined, grant2006disciplined}.
These include \texttt{CVX}~\cite{cvx} in MATLAB, \texttt{CVXPY}~\cite{diamond2016cvxpy, agrawal2018rewriting}
in Python, and \texttt{Convex.jl}~\cite{udell2014convex} in Julia.

\paragraph{Modern open-source modeling languages}
At the time of writing, most modeling languages mentioned above continue to be
developed and have active communities. In the remainder of the section, we focus
our attention on two in particular: \texttt{JuMP} for the design ideas introduced by
\texttt{MathOptInterface}, and
\texttt{CVXPY} because it is the system on which the present work builds---a choice
motivated by the author's previous experience as a \texttt{CVXPY} contributor.

\subsection{\texttt{JuMP} and \texttt{MathOptInterface}}

\texttt{JuMP} is an algebraic modeling language embedded in Julia, originally developed in 2012 by
Iain Dunning, Joey Huchette, and Miles Lubin at the MIT Operations Research
Center~\cite{dunning2017jump, lubin2023jump10recent}. Its abstraction layer, \texttt{MathOptInterface}
(\texttt{MOI})~\cite{legat2020mathoptinterface}, rethinks how a modeling language communicates with solvers,
moving away from the traditional approach of reducing every problem to a single standard form.

\paragraph{Function-set pairs}
\texttt{MOI} represents each constraint as a \emph{function-set pair}: a mathematical function paired with
the set it must belong to. For example, a linear inequality $a^T x \leq b$ is encoded as
(\texttt{ScalarAffineFunction}, \texttt{LessThan(b)}), while a second-order cone
constraint $\|Ax + b\|_2 \leq c^T x + d$ uses (\texttt{VectorAffineFunction},
\texttt{SecondOrderCone}). This is more granular
than committing to a single conic standard form as in~\eqref{eq:conic}---\texttt{MOI} preserves the
semantic type of each constraint. Solvers that accept specific constraint types can receive them
directly, without unnecessary reformulations.

\paragraph{Bridges}
When a solver does not natively support a particular function-set pair, \texttt{MOI} uses
\emph{bridges}---small, composable transformations between equivalent constraint
representations.
For instance, a bridge can split an interval constraint into separate less-than and greater-than
constraints. The bridge system maintains a directed hyper-graph of all known transformations and
finds the shortest path to the target solver's supported types. Adding a reformulation once in \texttt{MOI}
benefits all solvers, rather than requiring each solver wrapper to implement its own conversions.

\paragraph{Incremental problem construction}
Some solvers, notably \texttt{MOSEK}~\cite{mosekSolver} and
\texttt{HiGHS}~\cite{huangfu2018highs}, expose APIs that add constraints one at a
time. \texttt{MOI} supports this incremental mode natively: rather than assembling a
complete coefficient matrix and passing it to the solver in one call, \texttt{MOI}
can issue individual \texttt{add\_constraint} calls as the model is built.
For example, adding a cut $a^T x \leq b$ after an initial solve requires
only one such call, not a full reconstruction of the problem data. This is
particularly useful in algorithms that iteratively refine a formulation,
such as cutting-plane methods or branch-and-cut, where constraints are
added between successive solves.

\paragraph{Extensibility}
\texttt{MOI} makes it straightforward to support new constraint types. A solver declares which function-set
pairs it handles, and new bridges can be registered at any time. \texttt{Hypatia} \cite{coey2022solving}
\cite{ocpb:16, odonoghue:21, cederbergspectralmatrixcones}.

\subsection{\texttt{CVXPY}}

\texttt{CVXPY} is a domain-specific language for convex optimization embedded in Python,
that was originally developed by Steven Diamond and Stephen Boyd at Stanford
University~\cite{diamond2016cvxpy}. Akshay Agrawal later redesigned the
rewriting system in \texttt{CVXPY}~\cite{agrawal2018rewriting}, 
which enabled support for other disciplined
languages~\cite{agrawal2019dgp, agrawal2020dqcp, agrawal2019dpp}, and the
project has since grown through contributions from many others. Unlike \texttt{MOI},
which preserves constraint diversity, \texttt{CVXPY} channels all problems through a
single compilation pipeline targeting one unified standard form.

\paragraph{Unified standard form}
\texttt{CVXPY} rewrites every problem into a conic standard form with a quadratic objective:
\begin{equation}
\begin{array}{ll}
\text{minimize}   & \frac{1}{2}\, x^T P\, x + c^T x \\
\text{subject to} & Ax + b \in \mathcal{K},
\end{array}
\label{eq:cvxpy-standard}
\end{equation}
Because the rewriting system targets a unique form, solver interfaces only need to map from this
single representation. This centralization simplifies the codebase: one set of rewriting rules
rather than a combinatorial number of bridges, and every solver benefits from the same
transformation
pipeline.

\paragraph{Disciplined convex programming}
\texttt{CVXPY} implements DCP rules that automatically verify convexity and transform high-level expressions
into conic form. A user can write \texttt{cp.sum\_squares(A @ x - b)} and \texttt{CVXPY} will verify
convexity, introduce auxiliary variables, and emit the equivalent second-order cone constraint
transparently. In contrast, \texttt{JuMP} requires users to formulate conic constraints directly (\eg
specifying a rotated second-order cone rather than a sum-of-squares objective), demanding deeper
conic optimization knowledge.

\paragraph{Parameters and efficient re-solving}
\texttt{CVXPY} provides first-class support for parameters via \texttt{cp.Parameter}. When a problem
contains parameters, \texttt{CVXPY} compiles the structure once---sparsity pattern, cone types, variable
ordering---and stores a mapping from parameter values to numerical coefficients. Subsequent solves
only require evaluating this mapping and calling the solver, skipping the full canonicalization. The
\texttt{cvxpygen} library~\cite{schaller2022cvxpygen} extends this by generating standalone C code embedding
the compiled structure, enabling deployment without a Python runtime.

\section{Canonicalization and Coefficient Extraction}

The previous sections described what modeling languages do at a high level:
they accept mathematical formulations and produce solver-ready data. This
section examines \emph{how} that translation works, using a concrete example
to illustrate the canonicalization pipelines of both \texttt{MOI} and \texttt{CVXPY}.

\subsection{A worked example: least-squares}

Consider the least-squares problem
\begin{equation}
\begin{array}{ll}
\text{minimize} & \|Ax - b\|_2^2,
\end{array}
\label{eq:least-squares}
\end{equation}
with data $A \in \reals^{2 \times 2}$ and $b \in \reals^2$. This problem is
small enough to trace every step of the canonicalization, yet it exercises the
key mechanisms of both systems.

\paragraph{\texttt{MOI}'s bridge-based path}
In \texttt{JuMP}, writing \texttt{@objective(model, Min, sum((A*x - b).\^{}2))} creates
a \emph{scalar quadratic function}. \texttt{MOI} represents the objective as a
function-set pair
(\texttt{ScalarQuadraticFunction}, \texttt{MinSense}). If the target solver
accepts quadratic objectives directly (\eg \texttt{MOSEK}, \texttt{Gurobi}), no further
transformation is needed---\texttt{MOI} passes the quadratic coefficients as-is.

If the solver only supports conic constraints, \texttt{MOI}'s bridge system activates.
A \emph{quadratic-to-SOC bridge} rewrites the quadratic objective into a
second-order cone constraint by introducing an epigraph variable $t$ and
enforcing $\|Ax - b\|_2 \leq t$, with the objective changed to
$\text{minimize } t$. Each bridge is a small, self-contained transformation;
\texttt{MOI} composes as many as needed to reach the solver's supported types.

\paragraph{\texttt{CVXPY}'s compilation path}
In \texttt{CVXPY}, the same problem is written as
\begin{verbatim}
x = cp.Variable(2)
prob = cp.Problem(cp.Minimize(cp.sum_squares(A @ x - b)))
\end{verbatim}

\texttt{CVXPY} processes this through several stages:

\begin{enumerate}
\item \textbf{Expression tree.} \texttt{CVXPY} builds a tree whose root is
  \texttt{sum\_squares}, with children \texttt{sub}($-$), which in turn has
  children \texttt{matmul}($A$, $x$) and constant $b$.

\item \textbf{DCP verification.} The \texttt{sum\_squares} atom is convex, and
  its argument $Ax - b$ is affine; by DCP composition rules, the overall
  expression is convex.

\item \textbf{Epigraph reformulation.} The \texttt{sum\_squares} atom's
  canonicalization method introduces an auxiliary variable $t \in \reals$ and
  rewrites the objective as $\text{minimize } t$ subject to
  $\|Ax - b\|_2^2 \leq t$, which is equivalent to
  $\left\|\begin{smallmatrix} Ax - b \\ \frac{1}{2}(t - 1)
  \end{smallmatrix}\right\|_2 \leq \frac{1}{2}(t + 1)$---a second-order cone
  constraint.

\item \textbf{Conic standard form.} \texttt{CVXPY} assembles the constraint into the
  form $\mathcal{A}z + \beta \in \mathcal{K}$ where $z = (x, t)$, and
  $\mathcal{K}$ is a second-order cone. The objective becomes $c^T z$ with
  $c = (0, 0, 1)$, selecting the epigraph variable.
\end{enumerate}

For our $2 \times 2$ example with $z = (x_1, x_2, t)^T$, the final
solver-ready data is:
\begin{equation}
  c = \begin{pmatrix} 0 \\ 0 \\ 1 \end{pmatrix}, \quad
  \mathcal{A} = \begin{pmatrix}
    0 & 0 & \frac{1}{2} \\[2pt]
    0 & 0 & -\frac{1}{2} \\[2pt]
    a_{11} & a_{12} & 0 \\
    a_{21} & a_{22} & 0
  \end{pmatrix}, \quad
  \beta = \begin{pmatrix}
    \frac{1}{2} \\ \frac{1}{2} \\ -b_1 \\ -b_2
  \end{pmatrix},
  \label{eq:ls-coefficients}
\end{equation}
with $\mathcal{K} = \mathcal{Q}^4$ (second-order cone of dimension 4). The
matrices $\mathcal{A}$, $\beta$, and $c$ are the \emph{coefficient data} that
this thesis is concerned with extracting.

\subsection{From atoms to coefficients}

The least-squares example illustrates a general pattern. Every \texttt{CVXPY} problem
goes through the same pipeline, regardless of complexity.

\paragraph{What is coefficient extraction?}
Coefficient extraction is the process of transforming a high-level convex
optimization formulation into the numerical arrays that a solver
consumes---the matrices $P$, $\mathcal{A}$, vectors $c$, $\beta$, and cone
specifications $\mathcal{K}$ of the standard
form~\eqref{eq:cvxpy-standard}. This is the final, computational step of the
canonicalization pipeline: once the problem has been rewritten into conic
form, the system must traverse the resulting expression trees and assemble
concrete numerical data.

\paragraph{The canonicalization pipeline}
\texttt{CVXPY}'s full pipeline proceeds in three phases:
\begin{enumerate}
\item \textbf{DCP verification.} Each node in the expression tree is checked
  against DCP composition rules. If the overall problem is verified as convex
  (or concave for maximization), canonicalization proceeds; otherwise, \texttt{CVXPY}
  raises an error.

\item \textbf{Atom canonicalization.} Each non-affine atom (\eg
  \texttt{sum\_squares}, \texttt{norm}, \texttt{log}) is replaced by its
  conic equivalent. This step introduces auxiliary variables and constraints,
  producing a reformulated problem where every expression is affine and every
  constraint has the form $\mathcal{A}z + \beta \in \mathcal{K}$.

\item \textbf{Coefficient assembly.} The system traverses the affine
  expression trees that remain after atom canonicalization, evaluating each
  linear operator node and assembling the results into sparse coefficient
  matrices. Chapter~\ref{sec:Theme1} details how these linear operators are
  represented and composed.
\end{enumerate}

\paragraph{Parametric compilation}
When a problem contains \texttt{cp.Parameter} objects, \texttt{CVXPY} can separate
structure from data. The sparsity pattern and cone types are determined once
at compilation time; only the numerical values in $\mathcal{A}$ and $\beta$
change when parameters are updated. This separation is what makes it possible
to extract the coefficient computation as a standalone routine---a mapping
from parameter values to solver-ready arrays---which is the core idea behind
the C extractor developed in this thesis.

\paragraph{Contrast with \texttt{MOI}}
\texttt{MOI}'s bridge-based architecture achieves a similar end---translating high-level
formulations into solver-ready data---but through a different mechanism.
Rather than compiling to a single standard form, \texttt{MOI} applies targeted
transformations that preserve constraint semantics as long as possible.
This design favors flexibility and incrementality (constraints can be added or
removed without recompiling the entire problem), while \texttt{CVXPY}'s design favors
a clean compilation model where the entire problem is reduced to one
uniform representation. The trade-off is architectural: \texttt{MOI}'s approach is
more modular, while \texttt{CVXPY}'s approach makes the coefficient extraction step
uniform and predictable, which is advantageous for code generation.

%%
%% SUMMARY
%%
\section{Summary}

The evolution from manual encoding to algebraic modeling languages reveals a consistent trend toward
raising the abstraction level between modeler and solver. \texttt{JuMP}'s \texttt{MathOptInterface} extends this
trajectory with function-set pairs, composable bridges, and incremental construction, offering a
flexible architecture that preserves constraint semantics all the way to the
solver. \texttt{CVXPY} takes a complementary path: by funneling all problems
through a single conic standard form and providing DCP verification and first-class parameters, it
enables a compilation model where the problem structure is determined once and only numerical data
changes between solves. Neither system, however,
provides a self-contained C implementation of the coefficient extraction pipeline. The ability to
compile a \texttt{CVXPY} problem into a lightweight, dependency-free C module that maps parameters to
solver-ready coefficient matrices is the contribution of this thesis.
